\paragraph{How many readings does assurance require?}
The minimum fixed budget is three on every block where two adaptive
readings suffice but two fixed readings do not: 168 original and 57 external
blocks, spanning all six systems. Thus two adaptive readings replace three
necessary and sufficient fixed readings on 225 blocks, a one-third reduction.
Each replacement saves one reading in exchange for an additional information
round: the three fixed readings can be chosen together.
The original count includes the 165 matched-half cases and three further
blocks; only the matched cases establish exact batch success one half.
Using adaptive two on these blocks and the minimum fixed batch elsewhere
gives the complete-panel reading totals in Table~\ref{tab:minimum-reading-budget}.

\begin{table}[htbp]\centering\small
\caption{Reserved optional scalar readings for guaranteed joint half-width
precision. Each plan succeeds for every completion in the declared exact
model. Counts exclude the two common retained endpoints. The adaptive plan
uses a fixed fallback where no two-reading guarantee is available.}
\label{tab:minimum-reading-budget}
\begin{tabular}{lrr}\toprule
Quantity & Original & External\\\midrule
Blocks & 819 & 279\\
Three fixed versus two adaptive & 168 & 57\\
Minimum fixed readings, summed & 2625 & 743\\
Adaptive with fixed fallback, summed & 2457 & 686\\
Reading reduction across full panel & 6.40\% & 7.67\%\\\bottomrule
\end{tabular}
\end{table}

We calculated these budgets from the frozen coefficients before using any
optional reading answers.
Enumeration of all 32 fixed subsets per block agrees with the linear scan.
A separately implemented shortest-path checker verifies minimality and
checks feasible endpoint witnesses and the adaptive interval covers.
Online Resource~1 supplies each concrete fallback batch, source and system
summaries, and the fixed-budget histogram. The capacity and error controls
test how additional information changes the acquisition requirement.
